% Distilled blueprint for Evans, Appendix C (Calculus Facts).

\section{Calculus Facts}
\label{sec:calculus-facts}

\subsection{Boundaries}

Let $U\subset\mathbb R^n$ be open and bounded, and let $k\in\{1,2,\ldots\}$.

\begin{definition}[$C^k$ boundary]
\label{def:ck-boundary-regularity}
\dcref{appC:C.1-def-1}
\group{Boundary Regularity}
The boundary $\partial U$ is $C^k$ if, for every $x^0\in\partial U$, there are $r>0$ and $\gamma\in C^k(\mathbb R^{n-1})$ such that, after relabeling and reorienting the coordinate axes if necessary,
\[
U\cap B(x^0,r)
=\{x\in B(x^0,r)\mid
x_n>\gamma(x_1,\ldots,x_{n-1})\}.
\]
The boundary is $C^\infty$ if it is $C^k$ for every positive integer $k$, and it is analytic if the local functions $\gamma$ are analytic.
\end{definition}

\begin{definition}[Outward normal and normal derivative]
\label{def:outward-normal-and-normal-derivative}
\uses{def:ck-boundary-regularity}
\dcref{appC:C.1-def-2}
\group{Boundary Regularity}

(i) If $\partial U$ is $C^1$, its outward-pointing unit normal field is
\[
\nu=(\nu^1,\ldots,\nu^n).
\]

(ii) For $u\in C^1(\overline U)$, the outward normal derivative is
\[
\frac{\partial u}{\partial\nu}:=\nu\cdot Du.
\]
\end{definition}

For $x^0\in\partial U$ and a local graph $x_n=\gamma(x_1,\ldots,x_{n-1})$, define
\[
\begin{cases}
y_i=x_i=:\Phi^i(x) & i=1,\ldots,n-1,\\
y_n=x_n-\gamma(x_1,\ldots,x_{n-1})=:\Phi^n(x),
\end{cases}
\qquad y=\mathbf\Phi(x),
\]
and
\[
\begin{cases}
x_i=y_i=:\Psi^i(y) & i=1,\ldots,n-1,\\
x_n=y_n+\gamma(y_1,\ldots,y_{n-1})=:\Psi^n(y),
\end{cases}
\qquad x=\mathbf\Psi(y).
\]
Then $\mathbf\Phi=\mathbf\Psi^{-1}$, the boundary is mapped locally to $\{y_n=0\}$ with $U$ on the side $y_n>0$, and
\[
\det D\mathbf\Phi=\det D\mathbf\Psi=1.
\]

\subsection{Gauss--Green Theorem}

Throughout this subsection, $U\subset\mathbb R^n$ is open and bounded and $\partial U$ is $C^1$.

\begin{theorem}[Gauss--Green Theorem]
\label{thm:gauss-green-theorem}
\uses{def:ck-boundary-regularity, def:outward-normal-and-normal-derivative}
\dcref{appC:thm-1}
\group{Gauss-Green Theorem}

If $u\in C^1(\overline U)$, then for $i=1,\ldots,n$,
\[
(1)\qquad
\int_Uu_{x_i}\,dx
=\int_{\partial U}u\nu^i\,dS.
\]
\end{theorem}

\begin{theorem}[Integration-by-parts formula]
\label{thm:integration-by-parts-formula}
\uses{thm:gauss-green-theorem}
\dcref{appC:thm-2}
\group{Gauss-Green Theorem}

If $u,v\in C^1(\overline U)$, then for $i=1,\ldots,n$,
\[
(2)\qquad
\int_Uu_{x_i}v\,dx
=-\int_Uuv_{x_i}\,dx
+\int_{\partial U}uv\nu^i\,dS.
\]
\end{theorem}
\begin{proof}
\sketch Apply \cref{thm:gauss-green-theorem} to the product $uv$ and expand $(uv)_{x_i}$.
\end{proof}

\begin{theorem}[Green's formulas]
\label{thm:greens-formulas}
\uses{thm:integration-by-parts-formula, def:outward-normal-and-normal-derivative}
\dcref{appC:thm-3}
\group{Gauss-Green Theorem}

If $u,v\in C^2(\overline U)$, then
\[
\text{(i)}\qquad
\int_U\Delta u\,dx
=\int_{\partial U}\frac{\partial u}{\partial\nu}\,dS,
\]
\[
\text{(ii)}\qquad
\int_UDv\cdot Du\,dx
=-\int_Uu\Delta v\,dx
+\int_{\partial U}u\frac{\partial v}{\partial\nu}\,dS,
\]
and
\[
\text{(iii)}\qquad
\int_U(u\Delta v-v\Delta u)\,dx
=\int_{\partial U}
\left(u\frac{\partial v}{\partial\nu}
-v\frac{\partial u}{\partial\nu}\right)\,dS.
\]
\end{theorem}
\begin{proof}
\sketch Apply (2) to $u_{x_i}$ with the second factor $1$ and sum over $i$ to obtain (i). Apply (2) componentwise to $v_{x_i}u$ and sum to obtain (ii). Interchange $u$ and $v$ in (ii) and subtract to obtain (iii).
\end{proof}

\subsection{Polar Coordinates and the Coarea Formula}

\begin{theorem}[Polar coordinates]
\label{thm:polar-coordinates}
\dcref{appC:thm-4}
\group{Coarea Formula}
(i) If $f:\mathbb R^n\to\mathbb R$ is continuous and summable, then for every $x_0\in\mathbb R^n$,
\[
\int_{\mathbb R^n}f\,dx
=\int_0^\infty
\left(\int_{\partial B(x_0,r)}f\,dS\right)\,dr.
\]

(ii) In particular, for every $r>0$,
\[
\frac d{dr}\left(\int_{B(x_0,r)}f\,dx\right)
=\int_{\partial B(x_0,r)}f\,dS.
\]
\end{theorem}

Theorem \ref{thm:polar-coordinates} is the radial special case of the coarea identity.

\begin{theorem}[Coarea formula]
\label{thm:coarea-formula}
\dcref{appC:thm-5}
\group{Coarea Formula}
Let $u:\mathbb R^n\to\mathbb R$ be Lipschitz continuous and assume that, for almost every $r\in\mathbb R$, the level set
\[
\{x\in\mathbb R^n\mid u(x)=r\}
\]
is a smooth $(n-1)$-dimensional hypersurface. If $f:\mathbb R^n\to\mathbb R$ is continuous and summable, then
\[
\int_{\mathbb R^n}f|Du|\,dx
=\int_{-\infty}^\infty
\left(\int_{\{u=r\}}f\,dS\right)\,dr.
\]
\end{theorem}

Theorem \ref{thm:polar-coordinates} follows from Theorem \ref{thm:coarea-formula} by taking $u(x)=|x-x_0|$; see \cite[Chapter 3]{EvansGariepy1992}.

\subsection{Convolution and Smoothing}

\begin{definition}[Notation: $U_\epsilon$]
\label{def:epsilon-interior-set-notation}
\dcref{appC:C.4-def-1}
\group{Mollification}
For open $U\subset\mathbb R^n$ and $\epsilon>0$, set
\[
U_\epsilon
:=\{x\in U\mid\dist(x,\partial U)>\epsilon\}.
\]
\end{definition}

\begin{definition}[Standard mollifier]
\label{def:standard-mollifier}
\dcref{appC:C.4-def-2}
\group{Mollification}
(i) Define $\eta\in C^\infty(\mathbb R^n)$ by
\[
\eta(x):=
\begin{cases}
C\exp\!\left(\dfrac{1}{|x|^2-1}\right) & |x|<1,\\
0 & |x|\geq1,
\end{cases}
\]
where $C>0$ is chosen so that
\[
\int_{\mathbb R^n}\eta\,dx=1.
\]

(ii) For $\epsilon>0$, set
\[
\eta_\epsilon(x)
:=\frac1{\epsilon^n}\eta\left(\frac x\epsilon\right).
\]
Then $\eta_\epsilon\in C^\infty(\mathbb R^n)$,
\[
\int_{\mathbb R^n}\eta_\epsilon\,dx=1,
\qquad
\spt(\eta_\epsilon)\subset B(0,\epsilon).
\]
The function $\eta$ is the \textit{standard mollifier}.
\end{definition}

\begin{definition}[Mollification]
\label{def:mollification}
\uses{def:standard-mollifier, def:epsilon-interior-set-notation}
\dcref{appC:C.4-def-3}
\group{Mollification}

If $f:U\to\mathbb R$ is locally integrable, its mollification on $U_\epsilon$ is
\[
f^\epsilon:=\eta_\epsilon*f,
\]
that is,
\[
f^\epsilon(x)
=\int_U\eta_\epsilon(x-y)f(y)\,dy
=\int_{B(0,\epsilon)}\eta_\epsilon(y)f(x-y)\,dy
\qquad(x\in U_\epsilon).
\]
\end{definition}

\begin{theorem}[Properties of mollifiers]
\label{thm:properties-of-mollifiers}
\uses{def:mollification, def:standard-mollifier}
\dcref{appC:thm-6}
\group{Mollification}

Let $f\in L^1_{\mathrm{loc}}(U)$. Then:

(i) $f^\epsilon\in C^\infty(U_\epsilon)$.

(ii) $f^\epsilon\to f$ almost everywhere as $\epsilon\to0$.

(iii) If $f\in C(U)$, then $f^\epsilon\to f$ uniformly on every compact subset of $U$.

(iv) If $1\leq p<\infty$ and $f\in L^p_{\mathrm{loc}}(U)$, then $f^\epsilon\to f$ in $L^p_{\mathrm{loc}}(U)$.
\end{theorem}
\begin{proof}
\sketch Differentiation under the integral is valid on $U_\epsilon$ because the kernel has compact support; for every multiindex $\alpha$,
\[
D^\alpha f^\epsilon(x)
=\int_UD^\alpha\eta_\epsilon(x-y)f(y)\,dy,
\]
which proves (i). At every Lebesgue point of $f$,
\[
(3)\qquad
\fint_{B(x,r)}|f(y)-f(x)|\,dy\longrightarrow0
\quad(r\to0),
\]
and $\eta_\epsilon\leq C\epsilon^{-n}$ on its support; hence
\[
|f^\epsilon(x)-f(x)|
\leq C\epsilon^{-n}
\int_{B(x,\epsilon)}|f(y)-f(x)|\,dy\longrightarrow0,
\]
proving (ii). For $V\subset\subset W\subset\subset U$, uniform continuity of $f$ on $W$ makes the same estimate uniform on $V$, proving (iii). Finally, for small $\epsilon$ and $h\in L^p(W)$, Jensen's inequality and Fubini give
\[
(4)\qquad
\|h^\epsilon\|_{L^p(V)}\leq\|h\|_{L^p(W)}.
\]
Choose $g\in C(W)$ with $\|f-g\|_{L^p(W)}<\delta$. Then
\[
\|f^\epsilon-f\|_{L^p(V)}
\leq2\delta+\|g^\epsilon-g\|_{L^p(V)},
\]
and (iii), followed by $\delta\downarrow0$, proves (iv).
\end{proof}

\subsection{Inverse Function Theorem}

Let $U\subset\mathbb R^n$ be open, let $\mathbf f=(f^1,\ldots,f^n):U\to\mathbb R^n$ be $C^1$, fix $x_0\in U$, and set $z_0=\mathbf f(x_0)$.

\begin{definition}[Notation: gradient matrix $D\mathbf f$]
\label{def:gradient-matrix-df-square-notation}
\uses{def:vector-valued-function-notation}
\dcref{appC:C.5-def-1}
\group{Inverse and Implicit Function Theorems}

The gradient matrix is
\[
D\mathbf f=
\begin{pmatrix}
f^1_{x_1}&\cdots&f^1_{x_n}\\
\vdots&\ddots&\vdots\\
f^n_{x_1}&\cdots&f^n_{x_n}
\end{pmatrix}_{n\times n}.
\]
\end{definition}

\begin{definition}[Jacobian $J\mathbf f$]
\label{def:jacobian-jf}
\uses{def:gradient-matrix-df-square-notation}
\dcref{appC:C.5-def-2}
\group{Inverse and Implicit Function Theorems}

The Jacobian is
\[
J\mathbf f
:=|\det D\mathbf f|
=\left|\frac{\partial(f^1,\ldots,f^n)}
{\partial(x_1,\ldots,x_n)}\right|.
\]
\end{definition}

\begin{theorem}[Inverse Function Theorem]
\label{thm:inverse-function-theorem}
\uses{def:jacobian-jf}
\dcref{appC:thm-7}
\group{Inverse and Implicit Function Theorems}

Assume $\mathbf f\in C^1(U;\mathbb R^n)$ and $J\mathbf f(x_0)\ne0$. Then there are open sets $V\subset U$ and $W\subset\mathbb R^n$, with $x_0\in V$ and $z_0\in W$, such that:

(i) $\mathbf f:V\to W$ is one-to-one and onto;

(ii) $\mathbf f^{-1}:W\to V$ is $C^1$;

(iii) if $\mathbf f\in C^k$ for $k=2,3,\ldots$, then $\mathbf f^{-1}\in C^k$.
\end{theorem}

\subsection{Implicit Function Theorem}

Let $n,m$ be positive integers.

\begin{definition}[Notation: points in $\mathbb R^{n+m}$]
\label{def:xy-point-notation-rnm}
\dcref{appC:C.6-def-1}
\group{Inverse and Implicit Function Theorems}
Write
\[
(x,y)=(x_1,\ldots,x_n,y_1,\ldots,y_m)
\]
for $x\in\mathbb R^n$ and $y\in\mathbb R^m$.
\end{definition}

Let $U\subset\mathbb R^{n+m}$ be open, let $\mathbf f=(f^1,\ldots,f^m):U\to\mathbb R^m$ be $C^1$, fix $(x_0,y_0)\in U$, and set $z_0=\mathbf f(x_0,y_0)$.

\begin{definition}[Notation: gradient matrix $D\mathbf f$]
\label{def:gradient-matrix-df-rectangular-notation}
\uses{def:xy-point-notation-rnm}
\dcref{appC:C.6-def-2}
\group{Inverse and Implicit Function Theorems}

The gradient matrix is
\[
D\mathbf f
=\begin{pmatrix}
f^1_{x_1}&\cdots&f^1_{x_n}&f^1_{y_1}&\cdots&f^1_{y_m}\\
\vdots&&\vdots&\vdots&&\vdots\\
f^m_{x_1}&\cdots&f^m_{x_n}&f^m_{y_1}&\cdots&f^m_{y_m}
\end{pmatrix}_{m\times(n+m)}
=(D_x\mathbf f,D_y\mathbf f).
\]
\end{definition}

\begin{definition}[Jacobian $J_y\mathbf f$]
\label{def:jacobian-jyf}
\uses{def:gradient-matrix-df-rectangular-notation}
\dcref{appC:C.6-def-3}
\group{Inverse and Implicit Function Theorems}

Set
\[
J_y\mathbf f
:=|\det D_y\mathbf f|
=\left|\frac{\partial(f^1,\ldots,f^m)}
{\partial(y_1,\ldots,y_m)}\right|.
\]
\end{definition}

\begin{theorem}[Implicit Function Theorem]
\label{thm:implicit-function-theorem}
\uses{def:jacobian-jyf}
\dcref{appC:thm-8}
\group{Inverse and Implicit Function Theorems}

Assume $\mathbf f\in C^1(U;\mathbb R^m)$ and $J_y\mathbf f(x_0,y_0)\ne0$. Then there are an open set $V\subset U$ containing $(x_0,y_0)$, an open set $W\subset\mathbb R^n$ containing $x_0$, and a $C^1$ map $\mathbf g:W\to\mathbb R^m$ such that:

(i) $\mathbf g(x_0)=y_0$;

(ii) $\mathbf f(x,\mathbf g(x))=z_0$ for every $x\in W$;

(iii) if $(x,y)\in V$ and $\mathbf f(x,y)=z_0$, then $y=\mathbf g(x)$;

(iv) if $\mathbf f\in C^k$ for $k=2,3,\ldots$, then $\mathbf g\in C^k$.

Thus $\mathbf g$ is locally determined by the equation $\mathbf f(x,y)=z_0$.
\end{theorem}

\subsection{Uniform Convergence}

\begin{theorem}[Arzel\`a--Ascoli Theorem]
\label{thm:arzela-ascoli-theorem}
\uses{def:uniform-equicontinuity}
\dcref{appC:thm-9}
\group{Uniform Convergence}

Let $\{f_k\}_{k=1}^\infty$ be real-valued functions on $\mathbb R^n$. Assume that some $M$ satisfies
\[
|f_k(x)|\leq M
\qquad(k=1,2,\ldots,\ x\in\mathbb R^n),
\]
and that the family is uniformly equicontinuous. Then there are a subsequence $\{f_{k_j}\}_{j=1}^\infty$ and a continuous function $f$ such that
\[
f_{k_j}\to f
\quad\text{uniformly on compact subsets of }\mathbb R^n.
\]
\end{theorem}

\begin{definition}[Uniform equicontinuity]
\label{def:uniform-equicontinuity}
\dcref{appC:def-9}
\group{Uniform Convergence}
The family $\{f_k\}_{k=1}^\infty$ is \textit{uniformly equicontinuous} if, for every $\varepsilon>0$, there exists $\delta>0$ such that
\[
|x-y|<\delta
\quad\Longrightarrow\quad
|f_k(x)-f_k(y)|<\varepsilon
\]
for all $x,y\in\mathbb R^n$ and $k=1,2,\ldots$.
\end{definition}
